\scanpage{64}
% archive foliation: 63
(i) Si $\underline W^a(\mathfrak X)$ et $\underline W^a(\mathfrak Y)$ sont définis, il en est de $=$ de $W^a(\mathfrak X\times_{W(S)}\mathfrak Y)$, et on a
\[
W^a(\mathfrak X\times_{W(S)}\mathfrak Y)\simeq
W^a(\mathfrak X)\times_S W^a(\mathfrak Y)
\]
et on a plus généralement \deleted{\unsure{pour}} les formules analogues pour les limites projectives finies [donc $(\alpha)$ l'objet unité $(\beta)$ les produits fibrés, d'où comme conséquences $(\gamma)$ Noyaux]

(ii)
\[
W^a(\mathfrak X)\times_S T\simeq W_T^a(\mathfrak X\times_{W(S)}W(T))
\]
[l'\unsure{existence} du premier mb implique celle du second; la définition de $W_T$ est claire --]

De (i) résulte que le foncteur $W^a$ transforme un préschéma en groupes, anneaux etc en préschéma de $=$ type.
